\documentclass[../main.tex]{subfiles}

\begin{document}

\chapter{度量空间: 一般性质}\label{chapter:ch09-度量空间}

在第 1 章我们建立了实数的三种类型的性质. 第一种类型的性质是与加法和乘法相关的代数性质. 第二种是正数的性质, 通过这种性质我们定义了序和绝对值的概念. 利用代数性质与序性质, 两个实数之间的距离定义为它们的差的绝对值. 实数具有的最后一个性质是完备性: 关于实数的完备性公理等价于实数集中的每个 Cauchy 序列收敛到某个实数. 在第 7 章开始的赋范线性空间的研究中, 实数的代数性质被推广到了线性空间; 绝对值被推广为范数的概念, 它诱导了点之间的距离的概念; 而实数的序性质被放到一边. 我们现在进一步推广概念. 本章的目标是研究被称为度量空间的一般空间, 对这种空间基本的概念是两点间的距离. 它没有线性结构. 欧氏空间中开集与闭集的概念就自然地推广到一般的度量空间, 序列的收敛以及函数或映射的连续性的概念也如此. 我们首先考虑这些一般性的概念, 接着研究具有较好结构的度量空间: 完备、紧或可分的度量空间.

\subfile{Chapter-09/section-01.tex}

\subfile{Chapter-09/section-02.tex}

\subfile{Chapter-09/section-03.tex}

\subfile{Chapter-09/section-04.tex}

\subfile{Chapter-09/section-05.tex}


\end{document}